\section{The Non-Hermitian Transfer Matrix}
\label{sec:nonhermitian}

One of the most important objects in the transverse picture is the spatial transfer matrix $\mathcal{E}$, a column of the rotated network. Unlike the Hamiltonian or the time-evolution operator $U(\delta t)$, this matrix is generally non-Hermitian---and, due to the asymmetry of the MPO (Section~\ref{sec:mporep}), it is also non-normal ($\mathcal{E}\mathcal{E}^\dagger\neq\mathcal{E}^\dagger\mathcal{E}$). Non-Hermitian operators behave differently from the Hermitian operators of standard quantum mechanics, and these differences are important both for the physics we want to extract and for the numerical methods we use below.

\subsection{Left and right eigenvectors, and biorthogonality}

A Hermitian operator has a single orthonormal eigenbasis. A non-Hermitian $\mathcal{E}$ has two: right eigenvectors, defined by $\mathcal{E}\ket{R_i}=\mu_i\ket{R_i}$, and left eigenvectors, by $\bra{L_i}\mathcal{E}=\mu_i\bra{L_i}$, and in general $\ket{L_i}\neq\ket{R_i}$. These sets are not individually orthonormal; instead they form a biorthogonal system, $\braket{L_i|R_j}\propto\delta_{ij}$. This has two direct consequences for the calculation. First, the dominant left and right temporal states must be obtained independently---the right state cannot be obtained by conjugating the left one as happens in the symmetric case. Second, the physically meaningful contraction is the bilinear overlap $\braket{L|R}$ that encodes the network, not the usual sesquilinear overlap $\langle L^*|R\rangle$. Therefore, overlaps must be taken without complex conjugation.

Since the eigenvalues $\mu_i$ are generally complex, we order them by their modulus. The eigenvalue with the largest modulus, $\mu_0$, controls the long-chain contraction and therefore determines the Loschmidt rate in the thermodynamic limit:
\begin{equation}
    \ell(T) = -\lim_{N\to\infty}\frac{1}{N}\log|\mathcal{L}(T)|
            = -\log|\mu_0(T)|.
    \label{eq:rate}
\end{equation}
For a finite open chain, this relation receives $\mathcal{O}(1/N)$ boundary corrections. Thus, when comparing $-\log|\mu_0|$ with a finite-size Loschmidt echo divided by $N$, these boundary effects have to be taken into account.

\subsection{Eigenvalues versus singular values}

For Hermitian positive operators, such as the RDMs used in conventional tensor-network methods, eigenvalues and singular values carry essentially the same information: the eigenvalues are real, positive, and coincide with the squared singular values. This is what makes the usual truncation procedure variationally controlled. For the non-Hermitian transfer matrix $\mathcal{E}$, and for the associated RTM, this connection is lost. Its trace, $\mathrm{Tr}\,\mathcal{E}=\sum_i\mu_i$, is a sum of complex eigenvalues, so different contributions can cancel even when the singular values are large. The canonical illustration is an object with eigenvalues $+1$ and $-1$: both moduli are large and would be preserved by an SVD-based truncation, yet they cancel exactly in the physical trace. As a result, truncating by singular values does not provide the same direct control over the physical contraction. The situation is further complicated by non-normality: finite applications of a non-normal matrix can show transient amplification not visible from the eigenvalues alone. In practice these effects remain manageable for the transfer matrices we study, but they explain why RTM truncation, although more efficient, is less rigorously controlled than RDM truncation (Section~\ref{sec:transverse:truncation}).

\subsection{Emergent dual unitarity as a spectral statement}
\label{sec:emergent}

Section~\ref{sec:cft:temporal} introduced emergent dual unitarity as a physical prediction of the underlying CFT: the transfer eigenvalues migrate onto the unit circle as $T\to\infty$, so that the temporal transfer matrix becomes effectively unitary. Here we restate it as a concrete statement about the spectrum of $\mathcal{E}$ that we can monitor numerically, and draw out its consequence for the power method. The important point is to separate two pieces of information carried by each eigenvalue: its modulus, which controls growth or decay under repeated contractions, and its phase, which carries the oscillatory dynamics.

For a unitary operator all eigenvalues lie on the unit circle, $\mu_j=e^{i\theta_j}$, so there is no hierarchy in modulus and the spectrum is distinguished only by its phases. This is true for the rows of our network, the time-evolution layers $U(\delta t)=e^{-iH\delta t}$, which are unitary by construction. The column transfer matrix $\mathcal{E}$ is different: at finite time it is not unitary, its eigenvalues move into the complex plane and develop a hierarchy in modulus, and the largest, $\mu_0$, dominates the contraction while subleading eigenvalues are suppressed by powers of $|\mu_1/\mu_0|<1$. This is why the Loschmidt rate is fixed by $\ell(T)=-\log|\mu_0(T)|$, and why the gap ratio $|\mu_1|/|\mu_0|$ also controls the convergence of the power method.

After a quench to a critical point, conformal invariance predicts that this hierarchy gradually disappears as $T$ grows. The leading eigenvalues move back toward a common circle: their moduli become degenerate while their phases remain distinct,
\begin{equation}
    \frac{|\mu_i(T)|}{|\mu_0(T)|}
        = 1 - \frac{a_i}{T} - \frac{b_i}{T^2} - \cdots
        \;\xrightarrow[T\to\infty]{}\; 1,
    \label{eq:gapclose}
\end{equation}
with corrections decaying as integer powers of $1/T$. A band of eigenvalues with equal moduli and distinct phases is the spectral signature of a unitary operator up to an overall scale---the sense in which $\mathcal{E}$ becomes effectively unitary at long times. Importantly, this comes from conformal invariance at the critical point, not from integrability. It also explains why we monitor the whole leading band rather than only $\mu_0$: once several eigenvalues are nearly degenerate in modulus, a single dominant eigenvalue is no longer clearly separated, a one-vector power method becomes unreliable, and one must track the leading band as a whole---which is exactly what the block power method of the next section is built to do.

\section{Resolving the Leading Spectrum: A Block Power Method}
\label{sec:blockpm}

The central methodological contribution of this work is an algorithm to resolve the leading spectrum of $\mathcal{E}$ through the gap closing that emergent dual unitarity forces. We motivate and state what it does here; the full construction---the oblique projection, the exact biorthogonal pairing, the Schur-stable subspace basis, and the two truncation modes---is developed from first principles in Appendix~\ref{app:blockpm}.

\subsection{Why a single dominant vector is not enough}

The power method of Section~\ref{sec:powermethod} converges to the dominant eigenpair at a rate set by the gap ratio $|\mu_1/\mu_0|$: each application of $\mathcal{E}$ suppresses the subleading components of the boundary states by one factor of $|\mu_1/\mu_0|<1$. But Equation~\eqref{eq:gapclose} says that a quench to a critical point drives exactly this ratio to one---emergent dual unitarity \emph{is} the statement that the leading moduli become degenerate as $T$ grows. The better the conformal physics, the worse the single-vector iteration: at the temporal extents where the interesting scaling sets in, it slows critically and eventually stalls, mixing the dominant eigenvector with its near-degenerate partners.

The cure, familiar from Hermitian eigensolvers, is to iterate a \emph{block} of $k$ vectors at once and let them span the whole quasi-degenerate cluster. Convergence of the leading pair is then governed not by $|\mu_1/\mu_0|$ but by the gap \emph{between the block and the rest of the spectrum}, which stays finite as long as the block is large enough to contain the cluster. The subtlety, and the reason a ready-made routine will not do, is that $\mathcal{E}$ is non-Hermitian and non-normal, so every ingredient of the block method---orthogonalization, projection, even the notion of ``overlap''---must be reformulated in the biorthogonal, bilinear language of Section~\ref{sec:nonhermitian}. That reformulation, implemented in the routine \texttt{block\_transfer\_eigs} we added on top of \texttt{ITransverse.jl}, is what Appendix~\ref{app:blockpm} carries out: one keeps independent left and right blocks, applies $\mathcal{E}$ and its bilinear transpose to them, and extracts the leading Ritz values through a small non-conjugating (Petrov--Galerkin) projection, de-mixing the block onto a well-conditioned Schur basis rather than onto the eigenvectors themselves.

\subsection{The \texorpdfstring{$\pm$}{+/-} partner and the physical gap}

One non-Hermitian feature must be stated in the main text, because every gap-closing plot depends on it. The transfer spectrum comes in near-$\pm$ pairs: for each dominant $\mu_0$ there is a partner of nearly equal modulus with phase shifted by $\approx\pi$, a period-two structure in the time direction. Two consequences follow. First, sorting by modulus lets $\mu_0$ and its partner swap places between iterations, which a naive difference would register as a spurious jump; the tracked values must be matched by continuity in the complex plane before differencing. Second, and more physically, the naive ratio $|\theta_2|/|\theta_1|$ measures the \emph{splitting of the pair}, not the spectral gap. The meaningful, physical gap uses the largest modulus among the Ritz values that are \textbf{neither} the physical $\mu_0$ (selected by continuity in $T$) \textbf{nor} its $-\mu_0$ partner:
\begin{equation}
  \mathrm{gap}(T) = \frac{\max\{|\theta_i| : i \neq i_0,\, i \neq i_{\text{partner}}\}}{|\theta_{i_0}|}.
  \label{eq:physgap}
\end{equation}
All gap-closing plots in this thesis use Equation~\eqref{eq:physgap}; the RTM- and RDM-truncated results agree on it, which excludes a truncation artifact. Along a ladder of temporal extents the converged block at one rung also \emph{warm-starts} the next, so the iteration only has to relax the newly added time sites---a practical acceleration that, as Section~\ref{sec:results} discusses, turns out to be load-bearing for tracking the correct physical branch as the gap tightens, not merely a speed-up.

\subsection{What bounds the reach}
\label{sec:reach}

The block method extends the usable window considerably---it is what carries the ANNNI-type analysis through the gap closing---but it does not abolish the wall of Section~\ref{sec:results:limits}, and the reason is worth stating precisely, because it separates what an algorithm can and cannot fix. The leading eigen\emph{values} of $\mathcal{E}$ are well-conditioned essentially always: their sensitivity involves the biorthogonal pair through $1/|\langle L_0|R_0\rangle|$, which the block method keeps under control. The individual eigen\emph{vectors}, by contrast, have condition number $\sim 1/\mathrm{gap}$: as emergent dual unitarity closes the gap, any numerical representation of ``the'' dominant eigenvector---and hence any entropy built from it---degrades at a rate fixed by the physics itself, not by the iteration. A block method with exact pairing and a Schur basis therefore delivers clean eigenvalues (moduli, phases, the gap~\eqref{eq:physgap}) arbitrarily deep into the dual-unitary regime, while the entropy profile, which needs the eigenvector, is trustworthy only in the pre-wall window where the leading pair is still isolated:
\begin{equation}
  \underbrace{\text{spectrum of } \mathcal{E}}_{\text{robust: eigenvalue conditioning}}
  \qquad\text{vs}\qquad
  \underbrace{S^{\mathrm{gen}}_2 \text{ from } (L_0, R_0)}_{\text{bounded by } 1/\text{gap}} .
  \label{eq:reach}
\end{equation}
Characterizing this boundary---where it sits, and why frustration pulls it to earlier times---is itself part of the results (Section~\ref{sec:results:limits}).

\subsection{Phase rigidity: a diagnostic for the wall}
\label{sec:rigidity}

The statement that the individual eigenvector has condition number $\sim1/\mathrm{gap}$ can be sharpened into a single measurable quantity, which we use in Section~\ref{sec:results:limits} to diagnose exactly \emph{how} the wall forms. The subtlety is genuinely non-Hermitian and has no counterpart in the Hermitian eigenproblems of standard tensor-network methods. For a Hermitian operator the left and right eigenvectors coincide and remain orthonormal; a small perturbation only rotates them gently, provided the eigenvalues stay apart. For our non-normal $\mathcal{E}$ the left and right eigenvectors are distinct objects paired by the bilinear form $\braket{L_i|R_j}$ (Section~\ref{sec:nonhermitian}), and they can do something a Hermitian pair cannot: as two eigenvalues approach each other, the corresponding left and right eigenvectors can \emph{coalesce}, so that the pair becomes self-orthogonal, $\braket{L|R}\to0$, even though neither vector vanishes. The limiting configuration is an \emph{exceptional point}, where the matrix is defective and can no longer be diagonalized at all.

The order parameter for this approach is the \emph{phase rigidity} of a biorthogonal pair~\cite{rotter2009},
\begin{equation}
    r_j = \frac{|\braket{L_j|R_j}|}{\lVert L_j\rVert\,\lVert R_j\rVert},
    \label{eq:rigidity}
\end{equation}
which equals one for a well-conditioned (near-Hermitian) pair and tends to zero as the pair approaches an exceptional point. Its usefulness is that it discriminates two scenarios that are otherwise easy to confuse. An ordinary avoided crossing keeps $r_j\approx1$ and nothing breaks; a collapse $r_j\to0$ signals that the eigenvectors themselves are becoming ill-defined, so that any observable built from an individual eigenvector---the temporal entropy in particular---loses meaning, even while the eigenvalue it belongs to stays perfectly accurate. This is exactly the asymmetry of Equation~\eqref{eq:reach}, now with a number attached to it. We defer the measurement of $r_j(T)$, which turns out to collapse geometrically across several orders of magnitude as the wall is approached, to Section~\ref{sec:results:limits}.
